% =====================================================================
%  Hoja de referencia de LaTeX (cheat sheet) — suite de documentos LaTeX
%  Nivel 1 (Empezar) · article apaisado a 3 columnas · consulta rápida.
%  Compilar: latexmk -pdf cheatsheet.tex
% =====================================================================
\documentclass[9pt,a4paper,landscape]{extarticle}

\usepackage[utf8]{inputenc}
\usepackage[T1]{fontenc}
\usepackage{lmodern}
\usepackage{microtype}
\usepackage[spanish,es-noshorthands]{babel}
\usepackage{amsmath}\usepackage{amssymb}
\usepackage[landscape,margin=1cm]{geometry}
\usepackage{multicol}
\usepackage{booktabs}
\usepackage{xcolor}
\usepackage{tcolorbox}
\usepackage{enumitem}
\usepackage{hyperref}
\hypersetup{unicode, pdftitle={Hoja de referencia de LaTeX},
            pdfauthor={Rodrigo Martín Vidal Galán}, hidelinks}

\setlength{\parindent}{0pt}
\setlength{\columnsep}{1.2em}
\definecolor{tit}{RGB}{20,60,120}

% Encabezado de bloque temático
\newtcolorbox{blk}[1]{colback=tit!4, colframe=tit!70, fonttitle=\bfseries\small,
  title=#1, boxsep=2pt, left=3pt, right=3pt, top=2pt, bottom=2pt,
  before skip=4pt, after skip=4pt}
% comando \cmd = nombre de comando en mono
\newcommand{\K}[1]{\texttt{\textbackslash #1}}
% fila código -> efecto en dos "columnas" via \hfill
\newcommand{\R}[2]{\texttt{#1}\dotfill{\footnotesize #2}\\}

\begin{document}
\begin{center}
  {\Large\bfseries Hoja de referencia de \LaTeX}\quad
  {\footnotesize · suite de documentos · consulta rápida}
\end{center}
\vspace{-2pt}

\begin{multicols}{3}

\begin{blk}{Esqueleto mínimo}
\texttt{\textbackslash documentclass\{article\}}\\
\texttt{\textbackslash usepackage[utf8]\{inputenc\}}\\
\texttt{\textbackslash begin\{document\}}\\
\hspace*{1em}\dots texto \dots\\
\texttt{\textbackslash end\{document\}}\\[2pt]
Compilar: \texttt{latexmk -pdf a.tex}
\end{blk}

\begin{blk}{Estructura}
\K{section\{\}} · \K{subsection\{\}}\\
\K{subsubsection\{\}} · \K{paragraph\{\}}\\
\K{chapter\{\}} (report/book)\\
\K{tableofcontents}\\
\K{label\{k\}} \dots{} \K{cref\{k\}}\\
\K{footnote\{\dots\}}
\end{blk}

\begin{blk}{Formato de texto}
\K{textbf\{negrita\}} · \K{textit\{cursiva\}}\\
\K{texttt\{monoespaciada\}}\\
\K{emph\{énfasis\}} · \K{underline\{\}}\\
\K{textsc\{Versalitas\}}\\
\K{large} \K{small} \K{footnotesize}\\
\K{enquote\{comillas\}} (csquotes)
\end{blk}

\begin{blk}{Matemática: modos}
En línea: \texttt{\$ f(x)=x\^{}2 \$}\\
Bloque: \texttt{\textbackslash[ \dots{} \textbackslash]}\\
Numerada: \texttt{equation}\\
Varias: \texttt{align} (con \texttt{\&} y \texttt{\textbackslash\textbackslash})\\
Cita ec.: \K{eqref\{k\}}
\end{blk}

\begin{blk}{Matemática: construcciones}
\K{frac\{a\}\{b\}} · \K{sqrt[n]\{x\}}\\
\texttt{x\^{}\{2\}} · \texttt{x\_\{i\}} (sup/sub)\\
\K{sum\_\{i=1\}\^{}\{n\}} · \K{prod} · \K{int}\\
\K{lim\_\{x\textbackslash to0\}} · \K{frac\{d\}\{dx\}}\\
\K{binom\{n\}\{k\}} · \K{cdot} · \K{times}\\
\K{left(} \dots \K{right)} {\footnotesize(auto-tamaño)}
\end{blk}

\begin{blk}{Símbolos frecuentes}
\K{alpha}\,$\alpha$ \K{beta}\,$\beta$ \K{pi}\,$\pi$ \K{theta}\,$\theta$\\
\K{infty}\,$\infty$ \K{partial}\,$\partial$ \K{nabla}\,$\nabla$\\
\K{leq}\,$\leq$ \K{geq}\,$\geq$ \K{neq}\,$\neq$ \K{approx}\,$\approx$\\
\K{to}\,$\to$ \K{Rightarrow}\,$\Rightarrow$ \K{iff}\,$\iff$\\
\K{in}\,$\in$ \K{subset}\,$\subset$ \K{cup}\,$\cup$ \K{cap}\,$\cap$\\
\K{mathbb\{R\}}\,$\mathbb{R}$ \K{forall}\,$\forall$ \K{exists}\,$\exists$
\end{blk}

\begin{blk}{Listas}
\texttt{\textbackslash begin\{itemize\}}\\
\hspace*{1em}\K{item} \dots\\
\texttt{\textbackslash end\{itemize\}}\\[2pt]
\texttt{enumerate} (numerada)\\
\texttt{description} con \K{item[etiq.]}\\
Opciones: \texttt{[noitemsep]}, \texttt{[label=(\textbackslash alph*)]} (enumitem)
\end{blk}

\begin{blk}{Tablas (booktabs)}
\texttt{\textbackslash begin\{tabular\}\{@\{\}lcr@\{\}\}}\\
\hspace*{0.5em}\K{toprule}\\
\hspace*{0.5em}A \& B \& C \texttt{\textbackslash\textbackslash}\\
\hspace*{0.5em}\K{midrule}\\
\hspace*{0.5em}1 \& 2 \& 3 \texttt{\textbackslash\textbackslash}\\
\hspace*{0.5em}\K{bottomrule}\\
\texttt{\textbackslash end\{tabular\}}\\[2pt]
{\footnotesize Sin líneas verticales. \texttt{tabularx} = ancho fijo.}
\end{blk}

\begin{blk}{Figuras}
\texttt{\textbackslash begin\{figure\}[htbp]}\\
\hspace*{0.5em}\K{centering}\\
\hspace*{0.5em}\K{includegraphics[width=.6\textbackslash textwidth]\{f\}}\\
\hspace*{0.5em}\K{caption\{\dots\}} \K{label\{fig:k\}}\\
\texttt{\textbackslash end\{figure\}}\\[2pt]
{\footnotesize Graficar funciones: \texttt{pgfplots} (\K{addplot}).}
\end{blk}

\begin{blk}{Espaciado y saltos}
\K{quad} \K{qquad} (espacio math)\\
\K{,} \K{;} \K{!} (finos en math)\\
\K{newpage} · \K{clearpage}\\
\K{vspace\{1cm\}} · \K{hspace\{\}}\\
\K{noindent} · \K{centering}\\
{\footnotesize Línea en blanco = párrafo nuevo}
\end{blk}

\begin{blk}{Preámbulo base (matemática)}
\K{usepackage[T1]\{fontenc\}}\\
\K{usepackage\{lmodern\}}\\
\K{usepackage[spanish,es-noshorthands]\{babel\}}\\
\K{usepackage\{mathtools,amssymb\}}\\
\K{usepackage\{microtype\}}\\
\K{usepackage\{hyperref\}}\\
{\footnotesize (ver doc \emph{Estándares} / \emph{Plantillas})}
\end{blk}

\begin{blk}{Errores típicos}
{\footnotesize
\texttt{??} en el PDF $\to$ compilá 2 veces / \texttt{latexmk}.\\[2pt]
\texttt{Missing \$} $\to$ símbolo math fuera de \texttt{\$\dots\$}.\\[2pt]
\texttt{Undefined control sequence} $\to$ comando mal escrito o falta \K{usepackage}.\\[2pt]
\texttt{File X.sty not found} $\to$ instalá el paquete.}
\end{blk}

\end{multicols}

\vfill
\begin{center}\footnotesize
Suite de documentos \LaTeX{} · Rodrigo Martín Vidal Galán ·
\texttt{texdoc <paquete>} abre el manual de cualquier paquete.
\end{center}

\end{document}
